\section{Methodology}

We describe our benchmarking pipeline, covering three components: (1) stratified data preparation, (2) model architectures with pretrained initialization, and (3) training protocols. All implementation details are aligned with our released codebase.

\subsection{Dataset Preparation}

Caltech-101~\citep{fei2004learning} contains 9{,}144 colour images across 102 categories (101 object classes plus a \textit{background} class). Tab.~\ref{tab:dataset_summary} summarises the key statistics. Images vary in resolution and aspect ratio; the class distribution is moderately imbalanced, with most classes containing 40--80 samples while a few exceed 400.

\begin{table}[h]
    \caption{\hl{Model architectures.} Head dim denotes the input dimension of the replaced classifier layer.}
    \label{tab:arch_summary}
    \setlength{\tabcolsep}{4pt}
    \small
    \begin{tabular}{l p{3.5cm} r}
    \toprule
    Model & Key idea & Head dim \\
    \midrule
    ResNet-50~\citep{he2016deep}          & Residual skip connections     & 2048 \\
    EfficientNet-B2~\citep{tan2019efficientnet} & Compound depth/width/res.\ scaling & 1408 \\
    ViT-B/16~\citep{dosovitskiy2020image}      & Patch-based self-attention    & 768  \\
    ConvNeXt-Tiny~\citep{liu2022convnet}   & Modernised ConvNet design     & 768  \\
    EVA-02-Small~\citep{fang2024eva}       & MIM pre-training + CLIP distil.\ & 384  \\
    \midrule
    SVM + ResNet-18~\citep{donahue2014decaf}   & Frozen CNN features + RBF-kernel & 512  \\
    \bottomrule
    \end{tabular}
    \end{table}

\noindent We include all 102 categories and apply a stratified split to preserve class proportions. For each class $c$ with $n_c$ images:

\bigskip
\par\noindent
\vspace*{1em}
\begin{align}
\eqnmarkbox[Emerald]{ntrain}{n_c^{\text{train}}} &= \max\!\Big(1,\;\left\lfloor \eqnmarkbox[OliveGreen]{rtrain}{0.70} \cdot \eqnmarkbox[NavyBlue]{nc}{n_c} \right\rfloor\Big), \quad
\eqnmarkbox[WildStrawberry]{nval}{n_c^{\text{val}}} = \left\lfloor \eqnmarkbox[BurntOrange]{rval}{0.15} \cdot n_c \right\rfloor, \nonumber\\
\eqnmarkbox[Plum]{ntest}{n_c^{\text{test}}} &= n_c - n_c^{\text{train}} - n_c^{\text{val}},
\label{eq:split}
\end{align}
\annotate[yshift=1.8em]{above,right}{ntrain}{\tiny Train count}
\annotate[yshift=1.8em]{above,left}{rtrain}{\tiny Train ratio}
\annotate[yshift=1.8em]{above,right}{nc}{\tiny Class size}
\annotate[yshift=1.8em]{above,right}{nval}{\tiny Val count}
\annotate[yshift=1.8em]{above,right}{rval}{\tiny Val ratio}
\annotate[yshift=-1em]{below,right}{ntest}{\tiny Test (remainder)}
\vspace*{1em}

\noindent The $\max(1, \cdot)$ operator guarantees at least one training sample per class. A fixed random seed ($s = 42$) ensures reproducibility. The resulting split yields 6{,}352 / 1{,}326 / 1{,}466 samples for train / validation / test.

\subsection{Model Architectures}

We benchmark six methods spanning three paradigms: convolutional networks, vision transformers, and classical machine learning. For all deep models, we load ImageNet-pretrained weights and replace the final classification head with a linear layer mapping to $C = 102$ classes. Tab.~\ref{tab:arch_summary} summarises the architectures.

\begin{table}[h]
\centering
\caption{\hl{Model architectures.} Head dim: input dimension of the replaced classifier.}
\label{tab:arch_summary}
\setlength{\tabcolsep}{3pt}
\footnotesize
\begin{tabular}{llr}
\toprule
Model & Key idea & Head dim \\
\midrule
ResNet-50 \citep{he2016deep}               & Residual connections          & 2048 \\
EfficientNet-B2 \citep{tan2019efficientnet} & Compound scaling              & 1408 \\
ViT-B/16 \citep{dosovitskiy2020image}      & Patch self-attention          & 768  \\
ConvNeXt-Tiny \citep{liu2022convnet}        & Modernised ConvNet            & 768  \\
EVA-02-Small \citep{fang2024eva}            & MIM + CLIP distillation       & 384  \\
\midrule
SVM + ResNet-18 \citep{donahue2014decaf}    & Frozen CNN feat.\ + RBF-SVM  & 512  \\
\bottomrule
\end{tabular}
\end{table}

\textbf{CNNs.} ResNet-50 introduces residual skip connections that enable training of deeper networks without degradation. EfficientNet-B2 uses compound scaling to jointly optimise depth, width, and input resolution for better parameter efficiency. ConvNeXt-Tiny modernises the standard ConvNet design with larger kernels, layer normalisation, and inverted bottlenecks, closing the gap with transformer-based models.

\textbf{Vision Transformers.} ViT-B/16 splits each image into $16{\times}16$ patches and processes them with multi-head self-attention, learning global dependencies from the first layer. EVA-02-Small extends this paradigm with masked image modelling (MIM) pre-training on ImageNet-22k and CLIP feature distillation, yielding stronger representations at a compact model size; it is loaded via the \texttt{timm} library.

\textbf{Classical baseline.} The SVM baseline extracts 512-dim features from a frozen \texttt{ResNet-18}, standardises them, and trains an RBF-kernel SVM ($C{=}1$, $\gamma{=}\text{scale}$) without back-propagation, serving as a reference for how much end-to-end fine-tuning improves over fixed-feature classification.

\subsection{Training Details}

\textbf{Optimisation.} All deep models use AdamW~\citep{loshchilov2017decoupled} ($\lambda{=}10^{-4}$) with differential learning rates: backbone $\eta_{\text{bb}}$ and head $\eta_{\text{head}}$. \texttt{CNNs} use $10^{-4}$/$10^{-3}$; transformers (\texttt{ViT}, \texttt{EVA-02}) use $2{\times}10^{-5}$/$2{\times}10^{-4}$. We warm up linearly for 5 epochs, then apply cosine annealing:
\begin{equation}
\eta_t = \eta_{\min} + \tfrac{1}{2}(\eta_{\max} - \eta_{\min})\!\left(1 + \cos \tfrac{(t - T_w)\pi}{T - T_w}\right),
\end{equation}
where $T{=}30$ and $T_w{=}5$. Early stopping halts training after 5 epochs without validation improvement.

\textbf{Regularisation and augmentation.} We use label smoothing~\citep{szegedy2016rethinking} ($\epsilon{=}0.1$), gradient clipping (max norm 1.0), and mixed-precision (FP16) on CUDA. Training augmentations include random crop ($224^2$ from $256^2$), horizontal flip, colour jitter, and random erasing~\citep{zhong2020random} ($p{=}0.25$); at inference we centre-crop $224^2$. All images use ImageNet normalisation.

\textbf{Loss.} We minimise cross-entropy with label smoothing:

\vspace{1em}
\begin{equation}
\label{eq:loss}
\mathcal{L} = -\sum_{k=1}^{\eqnmarkbox[NavyBlue]{K}{K}} \left[ \left(1-\eqnmarkbox[OliveGreen]{eps}{\epsilon}\right)\, \mathbb{1}[y=k] + \frac{\epsilon}{\eqnmarkbox[NavyBlue]{K2}{K}} \right] \log \eqnmarkbox[WildStrawberry]{pk}{p_k},
\end{equation}
\annotate[yshift=1.8em]{above,left}{K}{\tiny Num.\ classes}
\annotate[yshift=1.8em]{above,right}{eps}{\tiny Label smoothing}
\annotate[yshift=-1em]{below,right}{pk}{\tiny Predicted prob.}
\vspace{1em}

\noindent where $K{=}102$ classes, $\epsilon{=}0.1$, and $p_k$ is the softmax output. Batch size is 32.

\subsection{Evaluation Metrics}

We report top-1 and top-5 accuracy, macro- and weighted-average precision, recall, and F1-score, per-class accuracy, and row-normalised confusion matrices on the held-out test set.

\subsection{Ablation Studies}

To isolate individual design choices, we conduct four ablation experiments on EVA-02-Small or the SVM baseline, varying one factor at a time while keeping all other settings identical to the main experiments:

\begin{enumerate}[leftmargin=*, itemsep=2pt, topsep=4pt]
    \item \textbf{Input resolution}: $112^2$ vs.\ $224^2$ (multiples of the EVA-02 patch size $14{\times}14$).
    \item \textbf{Data augmentation}: with vs.\ without (random crop, horizontal flip, colour jitter, random erasing).
    \item \textbf{Optimizer}: AdamW vs.\ SGD with separately tuned learning rates.
    \item \textbf{Feature type} (SVM): learned CNN features (ResNet-18, 512-dim) vs.\ hand-crafted HOG descriptors (8{,}100-dim).
\end{enumerate}
